\chapter{Reproducibility, remaining validation, and conclusion}
\label{ch:conclusion}\label{sec:conclusion}
\chapterpromise{A final account of the established statistical results, the executed empirical results, the unresolved implementation work, and the scope of the generalized method.}

\section{Established statistical results}
\BayesBreak is a finite Bayesian calculation over ordered partitions. For regular
exponential-family segment models with proper conjugate priors, exact segment marginal likelihoods
and posterior moments are obtained from sufficient statistics. Under a factorized partition prior,
sum-product dynamic programming yields fixed-count marginal likelihoods, $p(k\mid y)$,
changepoint marginals, changepoint-event probabilities, and Bayes curves. A separate max-sum
recursion yields the joint MAP partition. For aligned conditionally independent sequences with a
common partition, sequence-specific segment marginal likelihoods multiply. The finite latent-group
criterion has monotone alternating updates under exact maximization of its Jensen minorizer. For
nonconjugate segment models, the posterior-odds bounds are conditional on a uniform error bound for
every admissible segment score.

\section{Interpretation of the executed results}
The archive contains executed synthetic recovery, boundary-calibration, observation-model,
latent-group, nonconjugate-approximation, and finite-range timing studies, together with four
real-data fits. The real-data analyses are descriptive because independent changepoint labels are
not archived. Comparator scores that use the \BayesBreak joint-MAP changepoints as the reference
measure agreement rather than external accuracy. The two incompatible computations remain in the
result history but are excluded from comparator or predictive conclusions.

\section{Implementation and validation still required}
\begin{figure}[H]
\centering
\resizebox{0.96\textwidth}{!}{\begin{tikzpicture}[node distance=7mm and 8mm]
\node[successbox,text width=25mm] (sem) {Statistical statements and notation corrected};
\node[neutralbox,right=of sem,text width=24mm] (code) {Planned implementation corrections};
\node[neutralbox,right=of code,text width=25mm] (tests) {Unit, numerical, and statistical tests};
\node[neutralbox,right=of tests,text width=24mm] (rerun) {Controlled recomputation with new result identifiers};
\node[primarybox,right=of rerun,text width=24mm] (criteria) {Pre-specified acceptance criteria};
\draw[arrPrimary] (sem)--(code);
\draw[arrPrimary] (code)--(tests);
\draw[arrPrimary] (tests)--(rerun);
\draw[arrPrimary] (rerun)--(criteria);
\node[successbox,below=10mm of criteria,text width=25mm] (release) {Release after scientific and numerical review};
\node[draw=Warning,fill=WarningPale,diamond,aspect=1.8,left=15mm of release,align=center,font=\scriptsize] (pivot) {Criteria satisfied?};
\draw[arrPrimary] (criteria)--(pivot);
\draw[arrPrimary] (pivot)--node[above,font=\scriptsize]{yes}(release);
\draw[arr,draw=Warning,bend right=24] (pivot.west) to node[below,font=\scriptsize]{no: stop, correct, or narrow the claim}(tests.south);
\end{tikzpicture}
}
\caption{Remaining implementation and validation stages. Corrections to source code and controlled
recomputations are distinct from the mathematical and editorial revision of this volume.}
\label{fig:implementation-gates}
\end{figure}

The highest-priority implementation tasks are: explicit representation of segment cohesion and
boundary hazard; correction of the final objective returned by the latent-group optimizer;
family-correct posterior prediction for Beta-valued observations; rejection or explicit modeling of
out-of-range prediction coordinates; the one-to-one changepoint matching rule; validation of
comparator coordinate axes; complete result metadata; and resolution of the 17 bounded-run tests
that neither passed nor failed within the previous time limit. Corrected comparator or predictive
calculations must receive new result identifiers and must not overwrite the archived outputs.

\section{Conclusion}
The main contribution of \BayesBreak is a generalized hierarchical Bayesian segmentation method
that combines exponential-family segment integration with exact posterior computation over
ordered partitions and extends that construction to irregular designs, multiple aligned sequences,
known groups, and finite latent groups. The generality lies in the observation-model independence of
the partition recursion once mathematically valid segment marginal likelihoods are supplied. The
hierarchical structure lies in explicit probability models for shared or group-specific
changepoints and in the separately stated latent-group optimization criterion.

This formulation is grounded in conjugate analysis, product-partition models, Bayesian
multiple-changepoint recursions, hierarchical modeling, and dynamic programming. Its limitations
are equally mathematical: exactness depends on finite segment integrals and local factorization;
shared-changepoint pooling depends on alignment and conditional independence; nonconjugate results
depend on numerical error control; and empirical conclusions are limited by the available
references, replications, and comparator definitions. Stating those conditions precisely preserves
the breadth of the generalized method without replacing statistical claims by branding language.
